%%%%%%%%%%%%%%%%%%%%%%%%%%%%%%%%%%%%%%%%%%%%%%%%%%%%%%%%%%%%%%%%%%%%%%%%%%%%%%%%
%% 
\cleardoubleoddpage
\chapter{Introduction}
\label{sec:introduction}


\section{Motivation}
\label{sec:introduction:motivation}

As malware-assisted cybercrime continues to rise, the need for reliable tools and structured procedures to understand and mitigate malicious software grows accordingly. 
Static code analysis plays a crucial role in this process: it enables analysts to examine the internal logic of a program without executing it, 
making it a safe and effective first step in incident response and threat intelligence.

However, static analysis tools differ significantly in their capabilities and design philosophies. Some are optimized for analyzing .NET or Java applications, 
others originated as hexadecimal editors that evolved through community-driven development, and some were created by major intelligence agencies with 
decades of reverse engineering expertise. Despite this diversity, these tools share a common characteristic: they are freely available and widely used, 
yet their practical differences are often not systematically documented.

This thesis is motivated by the need to clarify which tools are best suited for different analysis scenarios, enabling practitioners to make informed decisions 
and ultimately improving the quality and efficiency of malware investigations.



\section{Objectives and Approach}
\label{sec:introduction:objectives}

The primary objective of this thesis is to conduct a practice-oriented comparison of selected static analysis tools commonly used in reverse engineering and 
malware investigation. Specifically, the tools Ghidra, Binary Ninja, Cutter, ILSpy, and Bytecode Viewer are evaluated with respect to their suitability for 
analyzing benign and malicious software samples.

To achieve this, each tool is examined according to several criteria: installation and configuration effort, usability, supported file formats, and the quality 
of the generated disassembly or decompilation output. The evaluation includes both benign test programs written specifically for this work, allowing direct comparison 
between source code and analysis output, and a real-world malware sample obtained from MalwareBazaar.

All experiments are conducted in Debian 13 virtual machine using the VirualBox virtualization software. 
This ensures reproducibility, isolation, and safe handling of malicious programs throughout the evaluation process.



%\section{Outline}
%\label{sec:introduction:outline}
%
%This thesis is structured as follows.
%
%Chapter 2 introduces the fundamental concepts of reverse engineering and static analysis that are relevant for understanding the tools evaluated in this work.
%
%Chapter 3 presents the selected analysis tools and explains the reasoning behind their inclusion in the thesis.
%
%Chapter 4 describes the evaluation criteria that are used to compare the tools.
%
%Chapter 5 outlines the methodology used during testing, including the setup of the analysis environment and the selection of samples.
%
%Chapter 6 contains the practical analysis of both benign and malicious samples and documents the observations made during testing.
%
%Chapter 7 provides a comparison of the results and highlights the strengths and weaknesses of each tool.
%
%%%Chapter 8 opens a discussion over what the study has brought to light.
%
%Finally, Chapter 8 concludes the thesis and summarizes the main findings, provides recommendations, and is followed by suggestions for future work.


%% 
%%%%%%%%%%%%%%%%%%%%%%%%%%%%%%%%%%%%%%%%%%%%%%%%%%%%%%%%%%%%%%%%%%%%%%%%%%%%%%%%
